% !TEX encoding = UTF-8 Unicode
% 予備審査レジュメ(特研1審査会のレジュメ 2025年度 を元にしたテンプレート)
% ヘッダーは先輩(2025年度 7307)の「創造工学専攻 情報工学コース / 年度 特別研究 2 審査会」に合わせた
% ビルド:thesis/resume/ で latexmk resume.tex → build/resume.pdf
\documentclass[10pt,uplatex,dvipdfmx]{jsarticle}

% 描画・色関連
\usepackage[dvipdfmx]{graphicx}
\usepackage[dvipdfmx]{xcolor}
\usepackage{here}

% レイアウト調整関連
\usepackage{geometry}
\usepackage{setspace}
\usepackage{fancyhdr}
\usepackage{float}
\usepackage{multicol}

% 数式関連
\usepackage{amsmath}
\usepackage{amssymb}
\usepackage{mathtools}
\usepackage{bm}

% 表・リスト関連
\usepackage{booktabs}
\usepackage{enumerate}

% テキスト装飾関連
\usepackage{url}

% キャプション関連
\usepackage{caption}
\usepackage{subcaption}

\graphicspath{{figures/}{../figures/}{../../experiments/}}
\geometry{a4paper, left=20mm,right=20mm,top=20mm,bottom=20mm}

\definecolor{tmcitblue}{cmyk}{0.86,0.91,0,0.04}
\pagestyle{fancy}
\fancyhf{}
\lhead{\color{tmcitblue}\footnotesize 東京都立産業技術高等専門学校}
\rhead{\color{tmcitblue}\footnotesize 創造工学専攻 情報工学コース\\2026 年度特別研究 2 審査会}
\lfoot{\color{tmcitblue}\footnotesize Tokyo Metropolitan College of Industrial Technology}
\renewcommand{\headrulewidth}{0.0pt}
\setlength{\headheight}{20.2781pt}

\begin{document}
\setstretch{1.00}
{\centering
\Large{空間熱干渉を考慮した状態推定と動的熱マージンに基づくファン制御(仮題)}\\
\vspace{6pt}
\large{7303 鈴木黎映}\\  % TODO 今年度の番号を確認(特研1では 6303)
\large{指導教員 知念 賢一}
\par}
\begin{multicols*}{2}

\section{はじめに}
% 背景(エッジ・ラックの RasPi クラスタ、専用空調なし)→ 課題(スロットリング、ファンを常に回す無駄)→ 目的
(執筆中)

\section{関連研究}
% 熱 RC モデルと予測制御、不確かさを考慮したマージン、データセンタの熱再循環。比較表(○△×)にすると分かりやすい

\section{提案手法}
% 熱モデル(RC から導出)、外部熱影響 H、カルマンフィルタ、Δ秒先予測と σ、動的マージン M = ασ + δ、ファン制御則

\section{実験}
% 実験環境(ラック上段 3 台、PoE+ HAT ファン、室温計測)、同定実験、熱干渉の確認、予測精度

\section{結果と考察}
% 図表を 1〜2 個。うまくいかなかった点も書く

\section{おわりに}
% まとめと今後の課題(ファン制御の本実験は 12 月)

\small{
\begin{thebibliography}{99}
  \bibitem{shi_edge_2016} Weisong Shi, Jie Cao, Quan Zhang, Youhuizi Li, Lanyu Xu, ``Edge Computing: Vision and Challenges,'' IEEE Internet of Things Journal, Vol.3, No.5, pp.637--646, 2016.
\end{thebibliography}
}

\end{multicols*}
\end{document}
